\begin{figure}[!htbp]
  \centering
  \resizebox{\linewidth}{!}{%
  \begin{tikzpicture}[
      stage/.style={
        draw=blue!55!black,
        rounded corners=2pt,
        fill=blue!4,
        align=center,
        text width=3.0cm,
        minimum height=1.35cm,
        inner sep=5pt,
        font=\small
      },
      note/.style={
        draw=black!35,
        rounded corners=2pt,
        fill=black!3,
        align=center,
        text width=2.95cm,
        inner sep=4pt,
        font=\scriptsize
      },
      arr/.style={-{Latex[length=2.2mm]}, thick, draw=black!55}
    ]
    \node[stage] (line) {긴 전보선ㆍ전화선\\신호 왜곡};
    \node[stage, right=6mm of line] (complex) {복소 임피던스\\$V=IZ$};
    \node[stage, right=6mm of complex] (rlc) {RLC 회로\\저장과 소산};
    \node[stage, right=6mm of rlc] (msd) {질량-스프링-댐퍼\\기계 임피던스};
    \node[stage, right=6mm of msd] (robot) {로봇 끝단\\가상 임피던스};
    \node[stage, right=6mm of robot] (lab) {MuJoCo Lab\\로그와 plot};

    \draw[arr] (line) -- (complex);
    \draw[arr] (complex) -- (rlc);
    \draw[arr] (rlc) -- (msd);
    \draw[arr] (msd) -- (robot);
    \draw[arr] (robot) -- (lab);

    \node[note, below=5mm of line] {왜 저항만으로\\설명이 부족한가?};
    \node[note, below=5mm of complex] {크기와 위상을\\함께 본다};
    \node[note, below=5mm of rlc] {$R$: 소산\\$L,C$: 저장};
    \node[note, below=5mm of msd] {$b$: 소산\\$m,k$: 저장};
    \node[note, below=5mm of robot] {힘과 운동의 관계를\\어떻게 설계하는가?};
    \node[note, below=5mm of lab] {파라미터 변화가\\응답을 어떻게 바꾸는가?};
  \end{tikzpicture}%
  }
  \caption{이 튜토리얼이 따라가는 임피던스 개념의 이동 경로. 이 그림은 임피던스의 완전한 역사 연대기가 아니라, 긴 통신선의 신호 왜곡이라는 중요한 역사적 동기에서 출발해 에너지 저장과 소산, 기계적 임피던스, 로봇 끝단의 접촉 반응 설계까지 이어지는 학습 경로를 요약한 것이다.}
  \label{fig:impedance-roadmap}
\end{figure}
